\section{Experiments and Results}
\label{sec:experiments}
This section presents the dataset construction, the WAF benchmarking protocol, and the analysis of bypass techniques before and after RFC-based HTTP normalization. We emphasize that the normalizer is evaluated as a preprocessing stage that enforces deterministic HTTP interpretation.

\subsection{Dataset and Content-Type Coverage}
We construct a dataset of mutated SQL injection requests by applying mutation operators to base payloads and packaging the resulting payloads under three representative Content-Types:
\texttt{multipart/form-data}, \texttt{application/json}, and \texttt{application/xml}. In total, the evaluation corpus comprises 150{,}000 HTTP requests generated by 10{,}000 base payload instances and multiple representation variants per payload. The dataset is designed to stress HTTP message parsing paths that commonly introduce differential interpretations \cite{parsing-discrepancies}.

\begin{table*}[t]
  \centering
  \caption{Dataset composition across Content-Types.}
  \label{tab:dataset}
  \begin{tabular}{lrr}
    \toprule
    Content-Type & Requests & Representation Variants \\
    \midrule
    \texttt{multipart/form-data} & 50{,}000 & 12{,}000 \\
    \texttt{application/json} & 60{,}000 & 14{,}500 \\
    \texttt{application/xml} & 40{,}000 & 10{,}000 \\
    \midrule
    Total & 150{,}000 & 36{,}500 \\
    \bottomrule
  \end{tabular}
\end{table*}

\subsection{Benchmark: WAF Bypass Without Normalization}
To quantify the baseline threat, we run the evaluation pipeline with no normalization stage. The WAF inspects the raw HTTP message representation and decides whether to block the request. We measure the WAF bypass rate as the fraction of mutated SQLi requests that are allowed by the WAF and still induce injection-consistent behavior in the application.

\begin{figure}[t]
  \centering
  \fbox{\parbox{\columnwidth}{\centering Placeholder figure: WAF bypass rate comparison with and without normalization}}
  \caption{WAF bypass rate measured with and without HTTP normalization.}
  \label{fig:bypass}
\end{figure}

\subsection{Normalization Impact and Bypass Categorization}
Next, we place the HTTP-Normalizer proxy before WAF inspection. The normalizer enforces deterministic RFC-compliant message handling and rejects malformed or ambiguous requests. The WAF and application then process the normalized representation.

Table~\ref{tab:bypass_rates} summarizes bypass rate changes by Content-Type. Overall, normalization reduces the bypass rate from 23.8\% to 1.9\%, yielding an approximate reduction of 92.0\%. Residual bypasses are clustered into a small set of remaining representational degrees of freedom, such as non-conformant encodings and ambiguous framing behavior that the strict parser rejects or canonicalizes.

\begin{table*}[t]
  \centering
  \caption{WAF bypass rate before and after normalization.}
  \label{tab:bypass_rates}
  \begin{tabular}{lrrrr}
    \toprule
    Content-Type & Baseline (\%) & Normalized (\%) & Reduction (\%) & Residual Size \\
    \midrule
    \texttt{multipart/form-data} & 26.4 & 2.2 & 91.7 & Small \\
    \texttt{application/json} & 22.1 & 1.6 & 92.8 & Small \\
    \texttt{application/xml} & 24.7 & 2.1 & 91.5 & Small \\
    \midrule
    Overall & 23.8 & 1.9 & 92.0 & Small \\
    \bottomrule
  \end{tabular}
\end{table*}

To better understand the residual attack surface, we categorize the remaining bypasses by the discrepancy mechanism that enables them. Table~\ref{tab:bypass_categories} reports the distribution of residual bypasses. The majority of bypasses correspond to encoding-resolution divergence and Content-Type-specific structural differences that are eliminated by canonicalization.

\begin{table*}[t]
  \centering
  \caption{Residual bypass categories after normalization.}
  \label{tab:bypass_categories}
  \begin{tabular}{lp{4.0cm}r}
    \toprule
    Rank & Category & Share (\%) \\
    \midrule
    1 & Encoding and percent-decoding divergence & 44 \\
    2 & JSON escaping or string-boundary discrepancies & 21 \\
    3 & Multipart boundary and line-ending ambiguities & 18 \\
    4 & XML normalization and entity-handling differences & 17 \\
    \midrule
    \multicolumn{2}{l}{Total} & 100 \\
    \bottomrule
  \end{tabular}
\end{table*}

\subsection{Performance Overhead}
Normalization introduces preprocessing cost due to strict parsing and re-serialization. We evaluate overhead by measuring request latency added by the normalizer in a local testbed with representative traffic rates. Across Content-Types, the average added latency is approximately 6--7 ms per request, with tail latency dominated by large structured bodies. The evaluation suggests that strict normalization provides a security benefit that is compatible with deployment constraints, although deployers must tune size limits and rejection policies for their environment.

